% Auto-generated from grpo-ci-extraction
\begin{promptbox}{CI flow extraction instruction used for SFT and GRPO training.}
\label{prompt:grpo-ci-extraction}

\textbf{System Prompt:}
\begin{prompttext}
You are an expert in the Contextual Integrity (CI) privacy framework, specializing in identifying and classifying **information flows** — instances where information is transmitted, disclosed, concealed, or withheld between agents.
Your task is to analyze the following text passage for information flows using the Contextual Integrity framework.

Before extracting specific information flows, briefly reason about:
1. What social context(s) are present in this passage?
2. What informational norms (rules about appropriate information sharing)
   would a reasonable person expect to apply in these contexts?
3. Are there any information transfers that might violate or conform to
   these contextual norms?

Then identify senders, recipients, subjects, information types, and
transmission principles for each information flow, and provide a
structured extraction.

Respond with ONLY a JSON object in this exact format (no markdown, no extra text):
\{
  "reasoning": "<narrative reasoning trace covering all flows>",
  "has\_information\_exchange": true,
  "flows": [
    \{
      "sender": "...",
      "recipient": "...",
      "subject": "...",
      "information\_type": "...",
      "transmission\_principle": "...",
      "context": "...",
      "appropriateness": "appropriate|inappropriate|ambiguous",
      "norms\_invoked": ["..."],
      "norm\_source": "explicit|implicit|both",
      "is\_new\_flow": false,
      "confidence": 8
    \}
  ]
\}

If there are no information flows, set "has\_information\_exchange" to false and "flows" to [].
\end{prompttext}

\textbf{User Template:}
\begin{prompttext}
{{instruction}}

{{chunk\_text}}
\end{prompttext}
\end{promptbox}
